\subsection{Multiplication operators}

Let X be a locally compact space and \(\mu\) a positive measure on X. We consider the multiplication operators on \(\mathrm{L}^{2}(\mathrm{X},\mu)\) ; these are closed operators with dense domain (prop. 5 of IV, p. 232). For any \(\mu\)-measurable function g on X, we will denote by \(m_{g}\) the partial multiplication operator by g on \(\mathrm{L}^{2}(\mathrm{X},\mu)\) .

PROPOSITION 22. — Let \(g\) be a function \(\mu\)-measurable on X.

\begin{enumerate}
\def\labelenumi{\alph{enumi})}
\item
  The spectrum of \(m_{g}\) is the \(\mu\)-essential image S of \(g\) ;
\item
  Let \(\lambda \in \mathbf{C} - \mathrm{Sp}(m_g)\) . The resolvent \(\mathrm{R}(m_g, \lambda)\) is the multiplication operator \(m_h\) , where \(h\) is the function on X defined by \(h(x) = 0\) if \(g(x) = \lambda\) and \(h(x) = (\lambda - g(x))^{-1}\) otherwise.
\end{enumerate}

Let us prove that \(\mathbf{C} - \mathbf{S}\) is contained in the resolvent set of \(m_g\) . Let \(\lambda \in \mathbf{C} - \mathbf{S}\) . There exists an open neighborhood U of \(\lambda\) such that the subset \(\mathrm{Y} = g^{-1}(\mathrm{U})\) of X is locally \(\mu\) -negligible. The function \(k\) defined on X by \(k(x) = (\lambda - g(x))^{-1}\) if \(x \notin \mathrm{Y}\) and \(k(x) = 0\) if \(x \in \mathrm{Y}\) belongs to \(\mathcal{L}^{\infty}(\mathrm{X}, \mu)\) (Lemma 5 of IV, p. 184); the operator of multiplication by \(k\) is therefore an endomorphism of \(\mathrm{L}^{2}(\mathrm{X}, \mu)\) .

Since \(|gk| \leqslant 1 + |\lambda k|\) , we have

\[
| g k f | \leqslant | f | + | \lambda k f |
\]

for \(f \in \mathcal{L}^2(\mathrm{X},\mu)\) , which implies that the image of \(m_k\) is contained in the domain of \(m_g\) . Conversely, let \(f \in \mathcal{L}^2(\mathrm{X},\mu)\) whose class \(\widetilde{f}\) belongs to the domain of \(m_g\) . Outside the locally \(\mu\) -negligible set Y, we have \(f(x) = k(x)(\lambda - g(x))f(x)\) , hence \(\widetilde{f}\) is in the image of \(m_k\) . The same formula proves that \(\lambda\) belongs to the resolvent set of \(m_g\) and that \(\mathrm{R}(m_g,\lambda) = m_k\) . Since the set Y is locally \(\mu\) -negligible, the multiplication operator \(m_k\) coincides with the operator \(m_h\) described in assertion \(b\) ).

Let us prove conversely that \(\mathbf{C} - \mathrm{Sp}(m_g)\) is contained in \(\mathbf{C} - \mathrm{S}\) . Let \(\lambda \in \mathbf{C} - \mathrm{Sp}(m_g)\) . Let \(\mathrm{M} > \| \mathrm{R}(m_g, \lambda) \|\) be a real number. Let Y denote the set of \(x \in \mathrm{X}\) such that \(|\lambda - g(x)| < \mathrm{M}^{-1}\) . Let us prove that Y is locally \(\mu\) -negligible, which will imply that \(\lambda\) does not belong to S, and will conclude the proof.

Let \(K\) be a compact subset of \(X\). Let \(\varphi\) be the characteristic function of \(Y \cap K\); this is an element of \(\mathcal{L}^{2}(X,\mu)\), and we denote its class in \(\mathrm{L}^{2}(X,\mu)\) by \(\widetilde{\varphi}\). Let \(\psi\) be a function in \(\mathcal{L}^{2}(X,\mu)\) whose class in \(\mathrm{L}^{2}(X,\mu)\) is \(\mathrm{R}(m_g,\lambda)(\widetilde{\varphi})\). We have \(\mathrm{R}(m_g,\lambda)(\widetilde{\varphi}) \in \operatorname{dom}(m_g)\) and \((\lambda - m_g)(\mathrm{R}(m_g,\lambda)(\widetilde{\varphi})) = \widetilde{\varphi}\), hence \((\lambda - g(x))\psi(x) = 1\) for \(\mu\)-almost all \(x \in Y \cap K\). This implies

\[
\| \mathrm{R} (m _ {g}, \lambda) \| ^ {2} \| \widetilde {\varphi} \| ^ {2} \geqslant \| \mathrm{R} (m _ {g}, \lambda) (\widetilde {\varphi}) \| ^ {2} \geqslant \mathrm{M} ^ {2} \mu (\mathrm{Y} \cap \mathrm{K}) = \mathrm{M} ^ {2} \| \widetilde {\varphi} \| ^ {2}.
\]

In view of the choice of M, this means that \(\varphi\) is zero \(\mu\) -almost everywhere. Thus, the set Y \(\cap\) K is \(\mu\) -negligible. The set Y is therefore locally \(\mu\) -negligible (INT, IV, p. 172, § 5, n° 2, prop. 5).

PROPOSITION 23. --- Let g be a μ-measurable function on X. The adjoint of the multiplication operator \(m_{g}\) is \(m_{\overline{g}}\) .

For every integer \(n \geqslant 1\) , let \(\varphi_{n}\) be the characteristic function of the set of elements \(x \in X\) such that \(|g(x)| \leqslant n\) and let \(\widetilde{\varphi}_{n}\) be its class in \(\mathrm{L}^{2}(\mathrm{X}, \mu)\) . Let \(f \in \mathcal{L}^{2}(\mathrm{X}, \mu)\) whose class \(\widetilde{f}\) belongs to \(\operatorname{dom}(m_{g}^{*})\) and let \(\psi\) be a function whose class is \(m_{g}^{*}(\widetilde{f})\) .

For every \(h \in \mathcal{L}^2(\mathrm{X}, \mu)\) whose class \(\widetilde{h}\) belongs to \(\operatorname{dom}(m_g)\) , we also have \(\widetilde{h}\widetilde{\varphi}_n \in \operatorname{dom}(m_g)\) , and hence \(\langle \widetilde{f} | m_g(\widetilde{h}\widetilde{\varphi}_n) \rangle = \langle m_g^*(\widetilde{f}) | \widetilde{h}\widetilde{\varphi}_n \rangle\) . This entails the equality

\[
\int_ {\mathrm{X}} \overline {{(f \overline {{g}} - \psi)}} \varphi_ {n} h d \mu = 0.
\]

Since the domain of \(m_{g}\) is dense in \(\mathrm{L}^{2}(\mathrm{X},\mu)\) , we deduce that \((f\overline{g}-\psi)\varphi_{n}\) is zero \(\mu\) -almost everywhere. Since n is arbitrary, this means that \(m_{g}^{*}(\widetilde{f})\) is the class in \(\mathrm{L}^{2}(\mathrm{X},\mu)\) of \(f\overline{g}\) . In particular, since \(m_{g}^{*}(\widetilde{f}) \in \mathrm{L}^{2}(\mathrm{X},\mu)\) we conclude that f belongs to the domain of \(m_{\overline{g}}\) and that \(m_{g}^{*}(\widetilde{f}) = m_{\overline{g}}(\widetilde{f})\) .

The adjoint of \(m_{g}\) is therefore an extension of \(m_{\overline{g}}\) . Moreover, we have

\[
\langle f \mid m _ {g} (h) \rangle = \int_ {\mathrm{X}} \overline {{f}} \cdot (g h) d \mu = \int_ {\mathrm{X}} f \overline {{g}} h d \mu
\]

for all \(f \in \mathrm{L}^{2}(\mathrm{X}, \mu)\) and \(h \in \mathrm{dom}(m_{g})\) , which proves that the linear form \(h \mapsto \langle f | m_{g}(h) \rangle\) is continuous when \(m_{\overline{g}}(f) g\) belongs to \(\mathrm{L}^{2}(\mathrm{X}, \mu)\) . Consequently, the domain of \(m_{\overline{g}}\) is contained in that of \(m_{g}^{*}\) , which concludes the proof.

COROLLARY. --- Let g be a μ-measurable function on X. The multiplication operator \(m_{g}\) on \(\mathrm{L}^{2}(\mathrm{X},\mu)\) is self-adjoint (resp. positive) if and only if the function g is locally μ-almost everywhere real-valued (resp. locally μ-almost everywhere positive).

The first assertion follows from the proposition. If \(g\) is locally \(\mu\)-almost everywhere positive, we have \(\langle f|m_{g}(f)\rangle=\int_{\mathrm{X}}g|f|^{2}d\mu\geqslant0\) for all \(f\in\mathrm{L}^{2}(\mathrm{X},\mu)\) , so the partial operator \(m_{g}\) is positive.

Conversely, if \(m_{g}\) is positive, then its spectrum is contained in \(\mathbf{R}_{+}\) (prop. 17 of IV, p. 248); since it is the \(\mu\)-essential image of g (prop. 22, a)), this means that g is locally \(\mu\) -almost everywhere positive.

Lemma 6. --- Let \(g_{1}\) and \(g_{2}\) be functions \(\mu\)-measurable on X.

\begin{enumerate}
\def\labelenumi{\alph{enumi})}
\item
  The partial operator \(m_{g_{1}} + m_{g_{2}}\) is closable and its closure is the multiplication operator \(m_{g_{1}+g_{2}}\) ;
\item
  One has \(m_{g_{1}} \circ m_{g_{2}} \subset m_{g_{1}g_{2}}\) ;
\item
  Suppose that \(g_{1}\) is bounded. We then have \(m_{g_{2}} \circ m_{g_{1}} = m_{g_{1}g_{2}}\) . Moreover, the domain of \(m_{g_{2}}\) is contained in the domain of \(m_{g_{1}g_{2}}\) , and \(m_{g_{1}} \circ m_{g_{2}}\) is the reduction of \(m_{g_{1}g_{2}}\) to \(\text{dom}(m_{g_{2}})\) .
\end{enumerate}

It is elementary that \(m_{g_{1}+g_{2}}\) is an extension of \(m_{g_{1}} + m_{g_{2}}\) ; this latter operator is therefore closable, and \(\overline{m_{g_{1}} + m_{g_{2}}} \subset m_{g_{1}+g_{2}}\) .

Let \(f \in \mathcal{L}^2(\mathrm{X},\mu)\) such that the function \(h = (g_1 + g_2)f\) belongs to \(\mathcal{L}^2(\mathrm{X},\mu)\) . For every integer \(n \geqslant 1\) , denote by \(\mathrm{X}_n\) the set of \(x \in \mathrm{X}\) such that \(|g_1(x)| + |g_2(x)| \leqslant n\) and let \(\varphi_n\) be the characteristic function of \(\mathrm{X}_n\) . We have \(\varphi_n f \in \mathrm{dom}(m_{g_1} + m_{g_2})\) . Since \((\varphi_n f)(x)\) converges to \(f(x)\) for all \(x \in \mathrm{X}\) and since for every \(n \in \mathbf{N}\) we have \(|\varphi_n f| \leqslant |f|\) and \(|(g_1 + g_2)\varphi_n f| \leqslant |(g_1 + g_2)f| = |h|\) , with \(h \in \mathcal{L}^2(\mathrm{X},\mu)\) by hypothesis, Lebesgue's theorem (INT, IV, p. 137, § 3, n° 7, th. 6) implies that the sequence of pairs of classes in \(\mathrm{L}^{2}(\mathrm{X},\mu)\times\mathrm{L}^{2}(\mathrm{X},\mu)\) of \((\varphi_{n}f,(g_{1}+g_{2})\varphi_{n}f)\) which belong to the graph of \(m_{g_{1}}+m_{g_{2}}\) converges to the pair of classes of \((f,h)\) in \(\mathrm{L}^{2}(\mathrm{X},\mu)\) . The closure of \(m_{g_{1}}+m_{g_{2}}\) is therefore indeed equal to \(m_{g_{1}+g_{2}}\) .

It is elementary that \(m_{g_{1}} \circ m_{g_{2}} \subset m_{g_{1}g_{2}}\) . Suppose that \(g_{1}\) is bounded, so that \(m_{g_{1}}\) is a continuous linear map on \(\mathrm{L}^{2}(\mathrm{X},\mu)\) . Then the domain of \(m_{g_{2}} \circ m_{g_{1}}\) is the set of equivalence classes of functions \(f \in \mathcal{L}^{2}(\mathrm{X},\mu)\) such that \(g_{2}(g_{1}f) \in \mathcal{L}^{2}(\mathrm{X},\mu)\) , that is, the domain of \(m_{g_{1}g_{2}}\) . We therefore have \(m_{g_{2}} \circ m_{g_{1}} = m_{g_{1}g_{2}}\) .

It is also elementary that \(\operatorname{dom}(m_{g_{1}} \circ m_{g_{2}}) = \operatorname{dom}(m_{g_{2}})\) is contained in \(\operatorname{dom}(m_{g_{1}g_{2}})\) , and that the restriction of \(m_{g_{1}g_{2}}\) to this space is equal to \(m_{g_{1}} \circ m_{g_{2}}\) .

One should beware of believing that \(m_{g_{1}} \circ m_{g_{2}} = m_{g_{1}g_{2}}\) in general (exercise 10 of IV, p. 347). Nevertheless, we have the following partial result:

PROPOSITION 24. --- Let g be a μ-measurable function on X.

\begin{enumerate}
\def\labelenumi{\alph{enumi})}
\item
  One has \(m_{\overline{g}} \circ m_g = m_{|g|^2}\) ;
\item
  For all integers \(k, \ell \in \mathbf{N}\) , we have \(m_{g^k\overline{g}^\ell} = m_g^k m_{\overline{g}}^\ell\) .
\end{enumerate}

One has \(m_{\overline{g}} \circ m_g \subset m_{|g|^2}\) by Lemma 6. Conversely, from the inequality \(|g| \leqslant 1 + |g|^2\) we deduce that \(\text{dom}(m_{|g|^2}) = \text{dom}(m_{\overline{g}} \circ m_g)\) , whence the first assertion.

Let \(k, \ell \in \mathbf{N}\) . We have \(m_g^k m_{\overline{g}}^\ell \subset m_{g^k \overline{g}^\ell}\) (loc. cit.). The domain of \(m_g^k m_{\overline{g}}^\ell\) is the set of classes in \(\mathrm{L}^2(\mathrm{X}, \mu)\) of functions \(h \in \mathcal{L}^2(\mathrm{X}, \mu)\) such that \(|g|^j h\) belongs to \(\mathcal{L}^2(\mathrm{X}, \mu)\) for every integer \(j\) such that \(0 \leqslant j \leqslant k + \ell\) . The inequalities

\[
\left| g ^ {j} h \right| \leqslant | h | + \left| g ^ {k + \ell} h \right|,
\]

valid for \(0 \leqslant j \leqslant k + \ell\) , allow one to observe that \(\mathrm{dom}(m_g^k m_{\overline{g}}^\ell)\) is equal to \(\mathrm{dom}(m_{g^k\overline{g}^\ell})\) , whence assertion \(b\) ).

